\chapter{Multiple Integrals}

In Chapters~32 and~33 we used partial and directional derivatives to dissect multivariable
functions into local slopes and tangent planes.
Now we turn to the reverse operation: accumulating a function of several variables across
a two-dimensional region or a three-dimensional solid.

Back in Chapter~20 we built the single integral $\int_a^b f(x)\,\dd x$ out of one
stubbornly simple idea: chop the interval $[a,b]$ into tiny pieces, build a thin
rectangle over each piece, add up their areas, and watch the sum settle onto a
number as the pieces shrink. The picture was a curve over a line, and the answer
was the area trapped between them.

Now the domain grows a dimension. Instead of integrating over an interval on the
$x$-axis, we integrate over a \emph{region} in the $xy$-plane -- a disk, a
triangle, a blob bounded by two curves. The graph of our function $f(x,y)$ is no
longer a curve but a \emph{surface} floating above that region, and the thing we
add up is no longer area but \emph{volume}: the volume of the solid trapped
between the region on the floor and the surface overhead.

The underlying idea has not changed at all: chop the region into tiny pieces,
estimate the contribution on each piece, add them up, and take the limit. What is
new is how we actually compute that sum: how to break a two-dimensional integral
into two ordinary integrals done one variable at a time (\textbf{iteration}, via
Fubini's theorem), and how to switch to curved coordinates when the region calls
for them (\textbf{change of variables}, via the Jacobian determinant). Those two
tools are the heart of this chapter.

\begin{focusbox}
A single integral adds a function over an \emph{interval} and measures
\emph{area}. A double integral adds a function over a \emph{planar region} and
measures \emph{volume}. The definition is the same Riemann sum (an approximating
sum of tiny column volumes), just with little patches of area $\Delta A$ instead
of little widths $\Delta x$. Everything hard is bookkeeping: how to organize the
sum (Fubini) and how to bend the coordinate grid without losing track of area
(the Jacobian).
\end{focusbox}

\chapterline

\section{The double integral as volume and as accumulation}
\label{ch34:double}

\subsection{Two readings of the same sum}

When you met $\int_a^b f(x)\,\dd x$ you learned to read it two ways at once, and
both readings will carry over. The \emph{geometric} reading: it is the signed
area under the graph of $f$. The \emph{accumulation} reading: if $f(x)$ is a
density-per-unit-length, then the integral is the total stuff. A double integral
has exactly these same two faces.

\textbf{Geometric reading.} Let $R$ be a region in the $xy$-plane and let
$f(x,y)\ge 0$ be a function defined on it. The graph $z=f(x,y)$ is a surface
hovering above the plane. The double integral
\[
  \iint_R f(x,y)\,\dd A
\]
is the \emph{volume} of the solid that sits over $R$ and reaches up to the
surface. When $f$ dips below zero the contribution is negative, exactly as area
went negative for the single integral when the curve dipped below the axis.

\textbf{Accumulation reading.} Suppose a flat metal plate occupies the region
$R$, and $f(x,y)$ is its \emph{density} -- grams per square centimetre -- which
may vary from point to point. Then $\iint_R f\,\dd A$ is the total \emph{mass} of
the plate. Here there is no surface and no volume in sight; we are simply adding
up a quantity that is spread out over an area, weighting each tiny patch by how
much stuff sits on it. Population over a county, rainfall over a watershed,
charge over a sheet -- all the same integral.

\begin{intuitionbox}
Think of $\dd A$ as an \emph{infinitesimal tile} of the floor, and $f(x,y)$ as
``how tall the building is'' (volume reading) or ``how heavy this tile is''
(mass reading) at that spot. The integral lays down infinitely many infinitely
small tiles over the whole region, multiplies each by its height-or-weight, and
sums. The $\iint$ -- a double S -- is just two nested ``sum'' signs, one for each
direction you have to sweep across the floor.
\end{intuitionbox}

\subsection{The Riemann sum, from scratch}

Let us build the thing so the definition is not a mystery. To keep
the picture clean, start with a region that is a rectangle,
$R=[a,b]\times[c,d]$ -- all points $(x,y)$ with $a\le x\le b$ and $c\le y\le d$.

Cut $[a,b]$ into $m$ equal sub-intervals of width $\Delta x=(b-a)/m$, and cut
$[c,d]$ into $n$ equal sub-intervals of height $\Delta y=(d-c)/n$. The grid lines
slice $R$ into $m\cdot n$ little rectangles, each of area
\[
  \Delta A=\Delta x\,\Delta y .
\]
In the $(i,j)$-th little rectangle pick any sample point
$(x_i^\ast,\,y_j^\ast)$. Over that rectangle, build a thin rectangular
\emph{box} whose base is the little rectangle and whose height is the value of
the function there, $f(x_i^\ast,y_j^\ast)$. The volume of that one box is
height times base area,
\[
  f(x_i^\ast,y_j^\ast)\,\Delta A .
\]
Now add up all $m\cdot n$ boxes to obtain an approximation to the total
volume under the surface:
\[
  S_{m,n}=\sum_{i=1}^{m}\sum_{j=1}^{n} f(x_i^\ast,y_j^\ast)\,\Delta A .
\]
This double sum is the two-dimensional Riemann sum.

\begin{figure}[H]
\centering
\begin{tikzpicture}
% Left panel: 2D floor grid
\begin{scope}[shift={(0,0)}]
  \draw[->] (-0.3,0) -- (4.8,0) node[right, font=\scriptsize] {$x$};
  \draw[->] (0,-0.3) -- (0,3.8) node[above, font=\scriptsize] {$y$};
  \foreach \i in {0,1,2,3,4} {
    \draw[gray!45] (\i,0) -- (\i,3.2);
  }
  \foreach \j in {0,1,2,3} {
    \draw[gray!45] (0,\j) -- (4,\j);
  }
  \fill[areafill,opacity=0.6, rounded corners=1pt] (2,1) rectangle (3,2);
  \node[font=\footnotesize] at (2.5,1.5) {$\Delta A$};
  \draw[<->, >=stealth] (0,-0.35) -- (1,-0.35) node[midway,below=1pt, font=\tiny] {$\Delta x$};
  \draw[<->, >=stealth] (-0.35,0) -- (-0.35,1) node[midway,left=1pt, font=\tiny] {$\Delta y$};
  \node[font=\scriptsize] at (2.5,-0.75) {Base tile $\Delta A = \Delta x\,\Delta y$};
\end{scope}

% Right panel: 3D column rising to surface
\begin{scope}[shift={(7.2,0)}]
\begin{axis}[
  width=7.2cm, height=6.2cm,
  view={45}{25},
  xlabel={$x$}, ylabel={$y$}, zlabel={$z$},
  xmin=0, xmax=3.5, ymin=0, ymax=3.5, zmin=0, zmax=4.5,
  tick label style={font=\tiny},
  title style={font=\small\bfseries},
  title={3D Column Accumulation},
]
% Surface overhead
\addplot3[surf, opacity=0.20, colormap/viridis, domain=0:3.2, y domain=0:3.2, samples=16]
  {3.2 - 0.2*x^2 - 0.15*y^2};
% Floor patch (1,1) to (2,2)
\addplot3[fill=areafill, opacity=0.5] coordinates {
  (1,1,0) (2,1,0) (2,2,0) (1,2,0) (1,1,0)
};
% 3D Column walls
\addplot3[fill=sourceblue!40, draw=sourceblue!80, opacity=0.7] coordinates {
  (1,1,0) (2,1,0) (2,1,2.5) (1,1,2.5) (1,1,0)
};
\addplot3[fill=sourceblue!50, draw=sourceblue!80, opacity=0.7] coordinates {
  (2,1,0) (2,2,0) (2,2,2.5) (2,1,2.5) (2,1,0)
};
% Column top cap
\addplot3[fill=sourceblue!30, draw=sourceblue!80, opacity=0.85] coordinates {
  (1,1,2.5) (2,1,2.5) (2,2,2.5) (1,2,2.5) (1,1,2.5)
};
\node[sourceblue, font=\tiny\bfseries, anchor=south, fill=white, inner sep=1.2pt, rounded corners=1pt] at (axis cs:1.6,1.6,3.55) {$\Delta V = f(x_i^*,y_j^*)\,\Delta A$};
\draw[->, >=stealth, sourceblue, thin] (axis cs:1.6,1.6,3.45) -- (axis cs:1.5,1.5,2.58);
\end{axis}
\end{scope}
\end{tikzpicture}
\caption{The double Riemann sum approximating the volume under a surface.
Left: The domain rectangle $R$ is partitioned into small tiles of area $\Delta A = \Delta x\,\Delta y$.
Right: Over each tile, we erect a rectangular column of height $f(x_i^*, y_j^*)$.
The volume of the column is $\Delta V = f(x_i^*, y_j^*)\,\Delta A$. As the grid is refined,
the sum of column volumes converges to the total volume under the surface $z=f(x,y)$.}
\label{ch34:fig-grid}
\end{figure}

If, as we make the grid finer and finer ($m,n\to\infty$, so
$\Delta x,\Delta y\to 0$), this sum approaches one and the same number no matter
how we chose the sample points, we call that number the \textbf{double integral}.

\begin{definition}[Double integral over a rectangle]
\label{ch34:def-double}
Let $f$ be defined on $R=[a,b]\times[c,d]$. If the Riemann sums $S_{m,n}$ above
converge to a single limit $L$ as $m,n\to\infty$, independent of the choice of
sample points, then $f$ is \emph{integrable} on $R$ and we write
\[
  \iint_R f(x,y)\,\dd A=L .
\]
\end{definition}

As in the single-variable setting, every continuous function on a closed bounded
region (a compact set in $\R^2$) is integrable. Continuity ensures that the Riemann
sums converge to a unique, well-defined limit as the partition mesh tends to zero.

\begin{happeningbox}
What does the differential $\dd A$ represent geometrically? In the finite
approximating sum, one multiplies a width by a height to obtain a patch of area
$\Delta A=\Delta x\,\Delta y$, weighting $f$ by that area. In the limit, $\dd A$
denotes the infinitesimal area element. The symbol $\dd A$ takes different forms
depending on the coordinate system: $\dd x\,\dd y$ in Cartesian coordinates,
$r\,\dd r\,\dd\theta$ in polar coordinates, and $|\det J|\,\dd u\,\dd v$ under a
general transformation. Each expression reflects the area of an infinitesimal
coordinate cell under that choice of coordinates.
\end{happeningbox}

\subsection{Regions that are not rectangles}

Real regions are rarely rectangles. To integrate over a triangle or a disk, do
the obvious thing: pretend $f$ is zero outside the region. Enclose the region $R$
inside a big rectangle, define a new function that equals $f$ on $R$ and equals
$0$ everywhere else in the rectangle, and integrate \emph{that} over the
rectangle. Tiles that fall outside $R$ contribute nothing; tiles inside
contribute $f\,\Delta A$ as before. The boundary tiles, the ones the curve slices
through, are a vanishing fraction of the total as the grid shrinks, so they do
not corrupt the limit. This little trick lets Definition~\ref{ch34:def-double}
cover any reasonable region without new ideas.

\begin{focusbox}
Useful sanity checks you will use constantly:
\begin{itemize}
\item $\displaystyle\iint_R 1\,\dd A=\text{area of }R$. (Height $1$ everywhere,
so the ``volume'' is just area times height $=$ area.)
\item $\displaystyle\iint_R c\,f\,\dd A=c\iint_R f\,\dd A$ and the integral of a
sum is the sum of the integrals -- linearity, inherited straight from the sums.
\item If $R$ splits into two non-overlapping pieces $R_1$ and $R_2$, then
$\iint_R=\iint_{R_1}+\iint_{R_2}$. Cut the region, add the volumes.
\end{itemize}
\end{focusbox}

\begin{recoverybox}
\textbf{Pause and think.}
Before calculating anything by formulas, evaluate $\iint_R (x^3 + \sin y)\,\dd A$
over the square $R = [-1,1]\times[-1,1]$.

\emph{Answer:} Exactly $0$. The surface overhead has two components: $x^3$ is odd across
the $y$-axis (meaning for every positive value at $(x,y)$ there is a matching negative value
over $(-x,y)$), and $\sin y$ is odd across the $x$-axis. By symmetry, the positive
and negative volumes cancel out completely. Spotting symmetry before computing limits
is one of the most useful habits in multivariable calculus.
\end{recoverybox}

\section{Iterated integrals and Fubini's theorem}
\label{ch34:fubini}

A double sum $\sum_i\sum_j$ is intimidating only until you remember you are
allowed to do the two sums \emph{one at a time}. Fix a column, add up that
column, then add up the column-totals. The continuous analogue of that regrouping
is the foundational computational tool of this chapter: \textbf{Fubini's theorem}, which states that a double
integral can be evaluated as two nested single integrals.

\subsection{The idea: volume by slices}

Here is the picture that makes Fubini obvious. You already know one way to find a
volume from Chapter~26: slice the solid into thin slabs and add the slab volumes.

Take our solid under $z=f(x,y)$ over the rectangle $[a,b]\times[c,d]$. Slice it
with a plane $x=\text{const}$, perpendicular to the $x$-axis. The slice is a thin
slab. Its \emph{cross-sectional area} is the area under the curve
$z=f(x,y)$ as $y$ runs from $c$ to $d$ -- with $x$ held fixed -- which is just a
single integral,
\[
  A(x)=\int_c^d f(x,y)\,\dd y .
\]
The slab has thickness $\dd x$, so its volume is $A(x)\,\dd x$. Add up the slabs
from $x=a$ to $x=b$:
\[
  \text{Volume}=\int_a^b A(x)\,\dd x
             =\int_a^b\!\left(\int_c^d f(x,y)\,\dd y\right)\dd x .
\]
That nested expression is an \emph{iterated integral}. We just sliced the solid
the long way and added the slices.

\begin{figure}[H]
\centering
\begin{tikzpicture}[scale=0.95]
  \begin{axis}[
    view={40}{25},
    width=10cm, height=7.5cm,
    xlabel={$x$}, ylabel={$y$}, zlabel={$z$},
    domain=0:2, y domain=0:2,
    samples=18, samples y=18,
    xtick=\empty, ytick=\empty, ztick=\empty,
  ]
    \addplot3[surf, opacity=0.55, shader=interp]
      {2 + 0.6*x*y};
    \addplot3[fill=tangentred!25, draw=tangentred, thick, opacity=0.6] coordinates {
      (1,0,0) (1,2,0) (1,2,3.2) (1,0,2) (1,0,0)
    };
    \addplot3[draw=tangentred, very thick, samples=30, samples y=1, domain=0:2]
      ({1},{x},{2+0.6*1*x});
    \node[tangentred!90!black, font=\small, anchor=south west] at (axis cs:1,1,1.5) {$A(x)$};
  \end{axis}
\end{tikzpicture}
\caption{Slicing the solid at a fixed $x$ gives a slab whose face area is
$A(x)=\int_c^d f(x,y)\,\dd y$ (the red profile). Adding the slabs
$A(x)\,\dd x$ from $a$ to $b$ rebuilds the volume. That is Fubini.}
\label{ch34:fig-slice}
\end{figure}

But we could just as well have sliced the \emph{other} way, with planes
$y=\text{const}$, giving
\[
  \text{Volume}=\int_c^d\!\left(\int_a^b f(x,y)\,\dd x\right)\dd y .
\]
The solid does not care which way you slice it. Both iterated integrals equal the
same volume, hence equal each other. That equality -- and the fact that either
equals the genuine double integral $\iint_R f\,\dd A$ -- is the content of
Fubini.

\begin{theorem}[Fubini, for nice functions]
\label{ch34:thm-fubini}
If $f$ is continuous on the rectangle $R=[a,b]\times[c,d]$, then
\[
  \iint_R f(x,y)\,\dd A
  =\int_a^b\!\int_c^d f(x,y)\,\dd y\,\dd x
  =\int_c^d\!\int_a^b f(x,y)\,\dd x\,\dd y .
\]
In particular both orders of integration give the same answer, and both equal the
double integral.
\end{theorem}

\begin{whybox}
Why is the iterated integral allowed to replace the double sum at all? Because
addition does not care how you group it. In the Riemann sum
$\sum_i\sum_j f(x_i^\ast,y_j^\ast)\,\Delta x\,\Delta y$, factor the constant
$\Delta x$ out of the inner sum:
\[
  \sum_i\Big(\underbrace{\sum_j f(x_i^\ast,y_j^\ast)\,\Delta y}_{\approx\,A(x_i^\ast)}\Big)\Delta x .
\]
The inner sum is a Riemann sum in $y$ alone; as the grid refines it becomes
$\int_c^d f(x_i^\ast,y)\,\dd y=A(x_i^\ast)$. The outer sum is then a Riemann sum
of $A(x)$ in $x$, becoming $\int_a^b A(x)\,\dd x$. So Fubini's theorem represents
``do the $y$-sum first, then the $x$-sum,'' the freedom to regroup a sum, carried
through the limit. The one piece of fine print -- which is why Fubini is a
\emph{theorem} and not a triviality -- is checking that the limits can be taken
in this nested order; continuity (more generally, $\iint|f|\,\dd A<\infty$) is
what licenses it.
\end{whybox}

\begin{connectionbox}
\textbf{Fine print: when is order-swapping safe?}
Iterating integrals by slicing was standard practice for Newton, Euler, and the
Bernoullis in the 1600s and 1700s, who sliced solids and added slices without
worrying about order. But swapping integrals is not always legal: if a function
oscillates wildly toward infinity, the two orders can yield genuinely different
numbers. In 1907, Guido Fubini proved the exact condition: swapping the order
is guaranteed safe whenever the function is \emph{absolutely integrable}
($\iint_R |f|\,\dd A < \infty$). Because every continuous function on a closed, bounded
domain automatically satisfies this, you can swap integration orders throughout
this book with complete confidence.
\end{connectionbox}

\subsection{A first computation}

\begin{example}[A double integral over a rectangle]
\label{ch34:ex-rect}
Compute $\displaystyle\iint_R (x+2y)\,\dd A$ over $R=[0,2]\times[0,3]$.

Integrate in $y$ first, treating $x$ as a constant:
\[
  \int_0^3 (x+2y)\,\dd y
  =\eval{xy+y^2}{0}^{3}
  =3x+9 .
\]
Now integrate that in $x$:
\[
  \int_0^2 (3x+9)\,\dd x
  =\eval{\tfrac{3}{2}x^2+9x}{0}^{2}
  =6+18=24 .
\]
As a check, slice the other way -- $x$ first:
\[
  \int_0^2 (x+2y)\,\dd x=\eval{\tfrac12 x^2+2xy}{0}^{2}=2+4y,
  \qquad
  \int_0^3 (2+4y)\,\dd y=\eval{2y+2y^2}{0}^{3}=6+18=24 .
\]
Same number, as Fubini promised.
\end{example}

\begin{trapbox}
The bounds and the differentials must match up \emph{from the inside out}. In
$\int_a^b\!\int_c^d f\,\dd y\,\dd x$, the inner differential $\dd y$ pairs with
the inner limits $c,d$, and the outer $\dd x$ pairs with $a,b$. A frequent
disaster is writing the limits in the wrong order, so that $x$ ends up
``integrated'' over $y$'s interval. Read the integral inside-out, like nested
parentheses, and keep each variable with its own bounds.
\end{trapbox}

\subsection{Non-rectangular regions: variable bounds}

For a rectangle the inner limits are constants. For a curvy region they are
\emph{functions}, because the cross-section changes length as you sweep across.
There are two standard shapes.

A region is \emph{$y$-simple} (vertically simple) if it sits between two curves
$y=g_1(x)$ below and $y=g_2(x)$ above, for $x$ in some interval $[a,b]$. Slice it
vertically: at each fixed $x$, $y$ runs from the bottom curve to the top curve.
\[
  \iint_R f\,\dd A=\int_a^b\!\int_{g_1(x)}^{g_2(x)} f(x,y)\,\dd y\,\dd x .
\]
A region is \emph{$x$-simple} (horizontally simple) if it lies between
$x=h_1(y)$ on the left and $x=h_2(y)$ on the right, for $y$ in $[c,d]$:
\[
  \iint_R f\,\dd A=\int_c^d\!\int_{h_1(y)}^{h_2(y)} f(x,y)\,\dd x\,\dd y .
\]

\begin{figure}[H]
\centering
\begin{tikzpicture}[scale=1.2]
  % Left panel: y-simple (vertical slices)
  \begin{scope}[shift={(0,0)}]
    \draw[->] (-0.2,0) -- (2.3,0) node[right, font=\scriptsize] {$x$};
    \draw[->] (0,-0.2) -- (0,2.3) node[above, font=\scriptsize] {$y$};
    \draw[domain=0:1, smooth, variable=\x, curveblue, very thick]
      plot ({2*\x},{2*\x*\x});
    \draw[domain=0:1, smooth, variable=\x, targetgreen, very thick]
      plot ({2*\x},{2*\x});
    \fill[areafill,opacity=0.35]
      plot[domain=0:1,smooth,variable=\x] ({2*\x},{2*\x})
      -- plot[domain=1:0,smooth,variable=\x] ({2*\x},{2*\x*\x})
      -- cycle;
    \draw[tangentred, thick, ->] (1.2,0.72) -- (1.2,1.2);
    \node[tangentred, right=3pt, font=\tiny] at (1.2,0.95) {slice};
    \node[targetgreen, above left, font=\scriptsize] at (1.5,1.55) {$y=g_2(x)$};
    \node[curveblue, below right, font=\scriptsize] at (1.5,0.9) {$y=g_1(x)$};
    \node[below left=1pt, font=\scriptsize] at (0,0) {$a$};
    \draw[gray] (2,0.02) -- (2,-0.02) node[below, font=\scriptsize] {$b$};
    \node[below, font=\footnotesize] at (1.0,-0.45) {$y$-simple: vertical slices};
  \end{scope}

  % Right panel: x-simple (horizontal slices)
  \begin{scope}[shift={(4.4,0)}]
    \draw[->] (-0.2,0) -- (2.3,0) node[right, font=\scriptsize] {$x$};
    \draw[->] (0,-0.2) -- (0,2.3) node[above, font=\scriptsize] {$y$};
    \draw[domain=0:1, smooth, variable=\y, curveblue, very thick]
      plot ({2*\y*\y},{2*\y});
    \draw[domain=0:1, smooth, variable=\y, targetgreen, very thick]
      plot ({2*\y},{2*\y});
    \fill[areafill,opacity=0.35]
      plot[domain=0:1,smooth,variable=\y] ({2*\y},{2*\y})
      -- plot[domain=1:0,smooth,variable=\y] ({2*\y*\y},{2*\y})
      -- cycle;
    \draw[tangentred, thick, ->] (0.72,1.2) -- (1.2,1.2);
    \node[curveblue, above left, font=\scriptsize] at (0.5,1.6) {$x=h_1(y)$};
    \node[targetgreen, below right, font=\scriptsize] at (1.5,1.4) {$x=h_2(y)$};
    \node[below left=1pt, font=\scriptsize] at (0,0) {$c$};
    \draw[gray] (0.02,2) -- (-0.02,2) node[left, font=\scriptsize] {$d$};
    \node[below, font=\footnotesize] at (1.0,-0.45) {$x$-simple: horizontal slices};
  \end{scope}
\end{tikzpicture}
\caption{Comparing slice directions for non-rectangular regions.
Left: A $y$-simple region sliced vertically. For each fixed $x\in[a,b]$, the inner variable $y$
runs from the lower boundary $y=g_1(x)$ to the upper boundary $y=g_2(x)$.
Right: An $x$-simple region sliced horizontally. For each fixed $y\in[c,d]$, the inner variable $x$
runs from the left boundary $x=h_1(y)$ to the right boundary $x=h_2(y)$.}
\label{ch34:fig-ysimple}
\end{figure}

\begin{example}[A triangular region]
\label{ch34:ex-tri}
Integrate $f(x,y)=x$ over the triangle $R$ with vertices $(0,0)$, $(1,0)$,
$(1,1)$. The hypotenuse is the line $y=x$; the bottom is $y=0$; the right edge is
$x=1$. Treating $R$ as $y$-simple: for each $x\in[0,1]$, $y$ runs from $0$ up to
$x$.
\[
  \iint_R x\,\dd A=\int_0^1\!\int_0^x x\,\dd y\,\dd x
  =\int_0^1 x\cdot\eval{y}{0}^{x}\,\dd x
  =\int_0^1 x\cdot x\,\dd x
  =\int_0^1 x^2\,\dd x=\frac13 .
\]
\end{example}

\subsection{Changing the order of integration}

Changing the order of integration is one of the most effective tools in multivariable
calculus. Frequently one order involves an inner integral with no elementary antiderivative,
while reversing the order introduces a factor that makes integration straightforward.
The procedure is systematic: \textbf{draw the region}, determine the limits in the
reverse order, and evaluate.

\begin{example}[An integral you cannot do in the given order]
\label{ch34:ex-order}
Evaluate
\[
  I=\int_0^1\!\int_x^1 e^{y^2}\,\dd y\,\dd x .
\]
The inner integral $\int e^{y^2}\,\dd y$ has \emph{no} elementary
antiderivative -- you are stuck cold. So change the order.

First decode the region from the bounds as written: $x$ goes from $0$ to $1$,
and for each $x$, $y$ goes from $y=x$ up to $y=1$. That is the triangle bounded
by $y=x$, $y=1$, and $x=0$ -- the triangle with vertices $(0,0)$, $(0,1)$,
$(1,1)$.

Now describe the \emph{same} triangle as $x$-simple. Let $y$ run from $0$ to $1$;
for each $y$, $x$ runs from the left edge $x=0$ to the right edge $x=y$ (the line
$y=x$ solved for $x$). So
\[
  I=\int_0^1\!\int_0^y e^{y^2}\,\dd x\,\dd y .
\]
Now the inner integral is trivial -- $e^{y^2}$ is constant in $x$:
\[
  \int_0^y e^{y^2}\,\dd x=y\,e^{y^2},
  \qquad
  I=\int_0^1 y\,e^{y^2}\,\dd y .
\]
And \emph{that} we can do, with the substitution $u=y^2$, $\dd u=2y\,\dd y$:
\[
  I=\frac12\int_0^1 e^{u}\,\dd u=\frac12(e-1).
\]
The same region, swept the other direction, turned an impossible integral into a
one-line substitution.
\end{example}

\begin{figure}[H]
\centering
\begin{tikzpicture}[scale=1.3]
  % Left panel: Given order (y-simple)
  \begin{scope}[shift={(0,0)}]
    \draw[->] (-0.2,0) -- (2.2,0) node[right, font=\scriptsize] {$x$};
    \draw[->] (0,-0.2) -- (0,2.2) node[above, font=\scriptsize] {$y$};
    \fill[areafill,opacity=0.35] (0,0) -- (0,1.8) -- (1.8,1.8) -- cycle;
    \draw[thick] (0,0) -- (0,1.8) -- (1.8,1.8) -- cycle;
    \draw[dashed, gray] (1.8,0) -- (1.8,1.8);
    \node[below, font=\scriptsize] at (1.8,0) {$1$};
    \node[left, font=\scriptsize] at (0,1.8) {$1$};
    \node[below left, font=\scriptsize] at (0,0) {$0$};
    \draw[tangentred, very thick, ->] (0.9,0.9) -- (0.9,1.8);
    \node[tangentred, right, font=\tiny] at (0.9,1.35) {$\dd y$ slice};
    \node[above, font=\scriptsize] at (0.9,1.8) {$y=1$};
    \node[below right, font=\scriptsize] at (1.1,0.9) {$y=x$};
    \node[below, font=\footnotesize] at (0.9,-0.4) {Given: vertical slice ($y$ from $x$ to $1$)};
  \end{scope}

  % Right panel: Reversed order (x-simple)
  \begin{scope}[shift={(4.5,0)}]
    \draw[->] (-0.2,0) -- (2.2,0) node[right, font=\scriptsize] {$x$};
    \draw[->] (0,-0.2) -- (0,2.2) node[above, font=\scriptsize] {$y$};
    \fill[areafill,opacity=0.35] (0,0) -- (0,1.8) -- (1.8,1.8) -- cycle;
    \draw[thick] (0,0) -- (0,1.8) -- (1.8,1.8) -- cycle;
    \draw[dashed, gray] (1.8,0) -- (1.8,1.8);
    \node[below, font=\scriptsize] at (1.8,0) {$1$};
    \node[left, font=\scriptsize] at (0,1.8) {$1$};
    \node[below left, font=\scriptsize] at (0,0) {$0$};
    \draw[targetgreen!80!black, very thick, ->] (0,1.2) -- (1.2,1.2);
    \node[targetgreen!80!black, above, font=\tiny] at (0.6,1.2) {$\dd x$ slice};
    \node[above, font=\scriptsize] at (0.9,1.8) {$y=1$};
    \node[below right, font=\scriptsize] at (1.1,0.9) {$x=y$};
    \node[below, font=\footnotesize] at (0.9,-0.4) {Reversed: horizontal slice ($x$ from $0$ to $y$)};
  \end{scope}
\end{tikzpicture}
\caption{Reversing the order of integration for Example~\ref{ch34:ex-order}.
The triangular region is identical in both panels.
Left: Slicing vertically requires an inner integral $\int_x^1 e^{y^2}\,\dd y$, which has no elementary antiderivative.
Right: Slicing horizontally swaps the inner variable to $x$, where $e^{y^2}$ is constant across the slice.
The inner integral trivially produces $y\,e^{y^2}$, setting up an easy $u$-substitution.}
\label{ch34:fig-order}
\end{figure}

\begin{walkbox}
Take Example~\ref{ch34:ex-order} on your walk. The number $\tfrac12(e-1)$ never
needed the antiderivative of $e^{y^2}$ -- yet the inner integral as first written
\emph{was} that forbidden antiderivative. How can a quantity be perfectly
computable while one route to it is blocked? (Hint: the blocked route asks for
the antiderivative as a \emph{function of $x$}; the open route only ever needs a
\emph{definite} integral, a single number. The double integral is a number; only
our bookkeeping made it look like it needed a forbidden function.)
\end{walkbox}

\begin{trapbox}
When you flip the order, you must re-derive the limits from the \emph{geometry},
not by mechanically swapping the symbols. Swapping
$\int_0^1\!\int_x^1\,\dd y\,\dd x$ into $\int_0^1\!\int_x^1\,\dd x\,\dd y$ is
nonsense -- you cannot have $x$ in the limits of an $x$-integral. Sketch the
region every single time; let the picture tell you the new bounds.
\end{trapbox}

\section{Polar coordinates and the Gaussian integral}
\label{ch34:polar}

Many regions are awkward to describe in rectangular coordinates. For domains
with circular symmetry -- disks, annuli, and circular sectors -- the natural
geometric parameters are radial distance and angle: \textbf{polar coordinates}
\[
  x=r\cos\theta,\qquad y=r\sin\theta,
\]
where $r\ge 0$ is the distance from the origin and $\theta$ is the angle. A disk
of radius $a$ is simply $0\le r\le a$, $0\le\theta\le 2\pi$ -- a rectangle in the
$(r,\theta)$ world.

\subsection{Why the polar area element carries an extra factor of $r$}

Here is the subtlety that trips up everyone the first time. In rectangular
coordinates, $\dd A=\dd x\,\dd y$, because a little grid cell is a tiny rectangle
of width $\dd x$ and height $\dd y$. In polar coordinates, $\dd A$ is
\[
  \boxed{\;\dd A=r\,\dd r\,\dd\theta\;}
\]
\emph{not} $\dd r\,\dd\theta$. There is an extra factor of $r$, and the whole
chapter's payoff is understanding where it comes from -- because it is the first
appearance of the change-of-variables idea that will close the chapter.

\begin{happeningbox}
Picture the polar grid: circles of constant $r$ and rays of constant $\theta$.
The little cell between radius $r$ and $r+\dd r$, and between angle $\theta$ and
$\theta+\dd\theta$, is \emph{not} a rectangle of sides $\dd r$ and $\dd\theta$.
Its two ``radial'' sides do have length $\dd r$. But its two ``arc'' sides are
pieces of circles, and the length of a circular arc is radius times angle:
$r\,\dd\theta$. So the cell is (very nearly) a rectangle of sides $\dd r$ and
$r\,\dd\theta$, with area
\[
  \dd A=(\dd r)(r\,\dd\theta)=r\,\dd r\,\dd\theta .
\]
Out at a larger radius $r$, the same angular wedge $\dd\theta$ sweeps out a wider
arc, so the tiles grow larger as you move away from the origin -- and the factor
$r$ tracks that growth.
\end{happeningbox}

\begin{figure}[H]
\centering
\begin{tikzpicture}
% Left panel: anatomy of one polar cell
\begin{scope}[shift={(0,0)}]
  \foreach \a in {15,30,45,60} {
    \draw[gray!40] (0,0) -- (\a:3.8);
  }
  \foreach \r in {1.4,2.4,3.4} {
    \draw[gray!40] (12:\r) arc (12:65:\r);
  }
  % Highlighted wedge
  \fill[areafill,opacity=0.6, rounded corners=1pt]
    (30:2.4) -- (30:3.4) arc (30:45:3.4) -- (45:2.4) arc (45:30:2.4) -- cycle;
  % Radial thickness
  \draw[tangentred, very thick] (30:2.4) -- (30:3.4)
    node[midway, below right, font=\scriptsize] {$\dd r$};
  % Outer arc
  \draw[curveblue, very thick] (30:3.4) arc (30:45:3.4)
    node[midway, above right, font=\scriptsize] {arc $= r\,\dd\theta$};
  \draw[->] (0,0) -- (4.0,0) node[right, font=\scriptsize] {$x$};
  \draw[->] (0,0) -- (0,3.6) node[above, font=\scriptsize] {$y$};
  \draw (0.6,0) arc (0:30:0.6) node[midway, right=1pt, font=\tiny] {$\theta$};
  \node[font=\footnotesize, below] at (2.0,-0.4) {$\dd A \approx (\dd r)(r\,\dd\theta) = r\,\dd r\,\dd\theta$};
\end{scope}

% Right panel: inner cell vs outer cell comparison lane
\begin{scope}[shift={(7.6,0)}]
  \foreach \a in {20,50} {
    \draw[gray!40] (0,0) -- (\a:4.0);
  }
  \draw[gray!40] (15:1.2) arc (15:55:1.2);
  \draw[gray!40] (15:1.9) arc (15:55:1.9);
  \draw[gray!40] (15:2.9) arc (15:55:2.9);
  \draw[gray!40] (15:3.6) arc (15:55:3.6);
  
  % Inner cell at r1 = 1.2
  \fill[targetgreen, opacity=0.4]
    (20:1.2) -- (20:1.9) arc (20:50:1.9) -- (50:1.2) arc (50:20:1.2) -- cycle;
  \node[font=\tiny, targetgreen!80!black, align=center] at (35:1.55) {Inner cell\\(small $r$)};
  
  % Outer cell at r2 = 2.9 (same dr, same dtheta)
  \fill[quotientorange, opacity=0.4]
    (20:2.9) -- (20:3.6) arc (20:50:3.6) -- (50:2.9) arc (50:20:2.9) -- cycle;
  \node[font=\tiny, quotientorange!90!black, align=center] at (35:3.25) {Outer cell\\(larger $r$)};

  \node[font=\scriptsize, align=center, below] at (2.0,-0.4) {
    Same $\Delta r$ and $\Delta\theta$, but different area!\\
    Area scales linearly with radius $r$.
  };
\end{scope}
\end{tikzpicture}
\caption{Why the polar area element carries an extra factor of $r$.
Left: A polar tile is bounded by two rays separated by $\dd\theta$ and two circular arcs separated by $\dd r$.
Its circumferential width is the arc length $r\,\dd\theta$, yielding area $\dd A = r\,\dd r\,\dd\theta$.
Right: Comparing two cells of identical $\Delta r$ and $\Delta\theta$ at different radii demonstrates that
tiles naturally widen as you move outward. The factor $r$ accounts for this geometric expansion.}
\label{ch34:fig-polar}
\end{figure}

So the conversion rule is: replace $x\to r\cos\theta$, $y\to r\sin\theta$, and
$\dd x\,\dd y\to r\,\dd r\,\dd\theta$, then set the limits as $r$- and
$\theta$-ranges.
\[
  \iint_R f(x,y)\,\dd x\,\dd y
  =\iint_{R'} f(r\cos\theta,\,r\sin\theta)\;r\,\dd r\,\dd\theta .
\]

\begin{example}[Area of a disk]
\label{ch34:ex-disk}
Take $f=1$ over the disk of radius $a$. In polar:
\[
  \iint_R 1\,\dd A=\int_0^{2\pi}\!\int_0^a 1\cdot r\,\dd r\,\dd\theta
  =\int_0^{2\pi}\eval{\tfrac12 r^2}{0}^{a}\,\dd\theta
  =\int_0^{2\pi}\frac{a^2}{2}\,\dd\theta=\pi a^2 .
\]
This cleanly yields the area formula $\pi a^2$, with the factor of $r$ playing an essential role.
\end{example}

\begin{recoverybox}
\textbf{Pause and think.}
Suppose you need the area of the upper half-disk $R = \{(x,y) : x^2+y^2\le a^2,\; y\ge 0\}$.
What are the integration limits in polar coordinates, and what does the integral compute?

\emph{Answer:} Radius $r$ sweeps from $0$ to $a$, but angle $\theta$ only sweeps from $0$
to $\pi$ (the upper half-plane). The integral is
$\int_0^\pi\!\int_0^a r\,\dd r\,\dd\theta = \pi \cdot \tfrac12 a^2 = \tfrac12\pi a^2$, exactly
half the area of the complete circle. The geometry sets the angle limits instantly.
\end{recoverybox}

\subsection{The Gaussian integral: evaluating $\int_{-\infty}^\infty e^{-x^2}\,\dd x$}

We can now evaluate one of the most important integrals in analysis, probability,
and physics:
\[
  I=\int_{-\infty}^{\infty} e^{-x^2}\,\dd x .
\]
The function $e^{-x^2}$ has no elementary antiderivative, so standard single-variable
techniques cannot evaluate it directly. Yet the value is exactly $\sqrt{\pi}$.
The classical solution, introduced by Poisson and Laplace, proceeds by squaring
the integral, viewing $I^2$ as a double integral over $\R^2$, and converting to
polar coordinates.

\begin{example}[$\int_{-\infty}^\infty e^{-x^2}\,\dd x=\sqrt\pi$]
\label{ch34:ex-gauss}
We cannot find $I$ directly, so consider $I^2$. Because the integral is just a
number, we may name the dummy variable of the second copy differently -- call it
$y$:
\[
  I^2=\left(\int_{-\infty}^\infty e^{-x^2}\,\dd x\right)
      \left(\int_{-\infty}^\infty e^{-y^2}\,\dd y\right)
     =\int_{-\infty}^\infty\!\int_{-\infty}^\infty e^{-x^2}e^{-y^2}\,\dd x\,\dd y .
\]
The key algebraic property is the radial factorization $e^{-x^2}e^{-y^2}=e^{-(x^2+y^2)}$.
The integrand depends on $(x,y)$ only through $x^2+y^2=r^2$, exhibiting radial symmetry.
Converting to polar coordinates turns this into an elementary integral.
The region is the whole plane, so $r$ runs $0$ to $\infty$ and $\theta$ runs $0$
to $2\pi$. With $x^2+y^2=r^2$ and the mandatory $\dd A=r\,\dd r\,\dd\theta$:
\[
  I^2=\int_0^{2\pi}\!\int_0^\infty e^{-r^2}\,r\,\dd r\,\dd\theta .
\]
Here the area factor $r$ in $\dd A=r\,\dd r\,\dd\theta$ provides the exact factor
needed for substitution. Setting $u=r^2$, $\dd u=2r\,\dd r$:
\[
  \int_0^\infty e^{-r^2}\,r\,\dd r
  =\frac12\int_0^\infty e^{-u}\,\dd u
  =\frac12\,\eval{-e^{-u}}{0}^{\infty}=\frac12 .
\]
Therefore
\[
  I^2=\int_0^{2\pi}\frac12\,\dd\theta=\frac12\cdot 2\pi=\pi,
  \qquad\Longrightarrow\qquad
  I=\sqrt\pi .
\]
\end{example}

\begin{aestheticbox}
A one-dimensional integral with no elementary antiderivative was evaluated by
passing to two dimensions, where radial symmetry could be exploited. The polar
factor $r$ in the area element supplied precisely the factor needed to evaluate
the radial antiderivative. Lifting a problem to a higher-dimensional setting where
new geometric symmetries become available is a recurring and powerful technique
in mathematics.
\end{aestheticbox}

\begin{applicationbox}
The bell curve of statistics is $\frac{1}{\sqrt{2\pi}\,\sigma}
e^{-(x-\mu)^2/2\sigma^2}$. The strange-looking constant $\frac{1}{\sqrt{2\pi}}$
is there for exactly one reason: so the total probability integrates to $1$.
Substituting and using $I=\sqrt\pi$ is what forces that normalizing constant.
Every time a statistician, physicist, or machine-learning engineer writes down a
Gaussian, the $\sqrt{2\pi}$ they carry around is a direct descendant of the polar
computation you just did.
\end{applicationbox}

\section{Triple integrals}
\label{ch34:triple}

Add one more dimension and nothing conceptually changes. A \textbf{triple
integral}
\[
  \iiint_E f(x,y,z)\,\dd V
\]
chops a solid region $E\subset\R^3$ into tiny boxes of volume
$\dd V=\dd x\,\dd y\,\dd z$, weights each by $f$, and sums. Fubini still holds:
evaluate it as three nested single integrals, in any of the six orders, with the
inner limits possibly functions of the outer variables.

\subsection{What it measures}

If $f=1$, the triple integral is the \emph{volume} of $E$:
$\iiint_E 1\,\dd V=\text{Vol}(E)$. (Just as $\iint 1\,\dd A$ was area.) If $f$ is
a density -- mass per unit volume -- the integral is the total \emph{mass} of a
solid body. There is no ``graph'' picture this time, because the graph
$w=f(x,y,z)$ would live in four dimensions; the accumulation reading is the one
to hold onto. We are summing a quantity spread through a three-dimensional body.

\begin{example}[Volume of a tetrahedron]
\label{ch34:ex-tet}
Find the volume of the solid $E$ bounded by the coordinate planes and the plane
$x+y+z=1$ (the corner of the unit cube cut by a slanted plane). Sweep $z$ first:
for fixed $(x,y)$ in the base triangle, $z$ runs from the floor $z=0$ up to the
slanted roof $z=1-x-y$. The base triangle (where the roof is still above the
floor, i.e.\ $1-x-y\ge 0$) is $0\le y\le 1-x$, $0\le x\le 1$. So
\[
  \text{Vol}(E)=\int_0^1\!\int_0^{1-x}\!\int_0^{1-x-y} 1\,\dd z\,\dd y\,\dd x .
\]
Innermost: $\int_0^{1-x-y}\dd z=1-x-y$. Middle:
\[
  \int_0^{1-x}(1-x-y)\,\dd y
  =\eval{(1-x)y-\tfrac12 y^2}{0}^{1-x}
  =(1-x)^2-\tfrac12(1-x)^2=\tfrac12(1-x)^2 .
\]
Outermost, with $u=1-x$:
\[
  \int_0^1\tfrac12(1-x)^2\,\dd x=\frac12\int_0^1 u^2\,\dd u=\frac12\cdot\frac13=\frac16 .
\]
The volume is $\tfrac16$, the familiar volume of that corner tetrahedron.
\end{example}

\begin{labbox}
Confirm $\tfrac16$ a second way. Treat the tetrahedron as a stack of horizontal
triangular slices at height $z$. At height $z$, the cross-section is the triangle
$x+y\le 1-z$, $x,y\ge 0$, a right triangle of legs $1-z$, so of area
$\tfrac12(1-z)^2$. Then $\text{Vol}=\int_0^1\tfrac12(1-z)^2\,\dd z=\tfrac16$.
Two slicing orders, one answer -- Fubini in three dimensions, by hand.
\end{labbox}

\subsection{Cylindrical and spherical coordinates}

Just as planar circular regions motivate polar coordinates, solids with axial or
spherical symmetry are most naturally evaluated in \emph{cylindrical} or
\emph{spherical} coordinates, where the geometry of coordinate cells determines
the volume element.

\textbf{Cylindrical} coordinates are just polar in the $xy$-plane with $z$ left
alone: $x=r\cos\theta$, $y=r\sin\theta$, $z=z$. A cell is a polar cell
($r\,\dd r\,\dd\theta$ of floor area) times a height $\dd z$, so
\[
  \dd V=r\,\dd r\,\dd\theta\,\dd z .
\]
The same factor $r$ as before, for the same arc-length reason.

\textbf{Spherical} coordinates measure a point by its distance $\rho$ from the
origin, the angle $\phi$ down from the positive $z$-axis (the ``polar angle'' or
colatitude, measuring how far you tip from the vertical pole), and the angle
$\theta$ around the $z$-axis (the ``azimuth'', the horizontal compass direction):
\[
  x=\rho\sin\phi\cos\theta,\quad
  y=\rho\sin\phi\sin\theta,\quad
  z=\rho\cos\phi .
\]
A spherical cell has three edges: a radial edge of length $\dd\rho$; an edge
along the $\phi$ direction, which is an arc of a circle of radius $\rho$, of
length $\rho\,\dd\phi$; and an edge along the $\theta$ direction, which is an arc
of a circle of radius $\rho\sin\phi$ (the distance from the $z$-axis), of length
$\rho\sin\phi\,\dd\theta$. Multiply the three edge lengths:
\[
  \boxed{\;\dd V=\rho^2\sin\phi\,\dd\rho\,\dd\phi\,\dd\theta\;}
\]

\begin{figure}[H]
\centering
\begin{tikzpicture}[scale=1.1, x={(-0.6cm,-0.3cm)}, y={(0.95cm,0.15cm)}, z={(0cm,1.0cm)}]
  % Axes
  \draw[->, gray!70] (0,0,0) -- (3.2,0,0) node[above left, font=\scriptsize] {$x$};
  \draw[->, gray!70] (0,0,0) -- (0,3.8,0) node[right, font=\scriptsize] {$y$};
  \draw[->, gray!70] (0,0,0) -- (0,0,3.8) node[above, font=\scriptsize] {$z$};

  % Coordinates of the spherical cell corners:
  \coordinate (O) at (0,0,0);
  \coordinate (P) at (1.62, 1.13, 1.98);
  \coordinate (Pproj) at (1.62, 1.13, 0);
  \coordinate (Pr) at (2.14, 1.50, 2.62);
  \coordinate (Pphi) at (2.02, 1.41, 1.34);
  \coordinate (Ptheta) at (1.09, 1.66, 1.98);
  \coordinate (Prphi) at (2.67, 1.87, 1.77);
  \coordinate (Prtheta) at (1.44, 2.20, 2.62);
  \coordinate (Pphitheta) at (1.36, 2.07, 1.34);
  \coordinate (Pcorner) at (1.80, 2.74, 1.77);

  % Dashed projection lines from origin to P
  \draw[dashed, gray!60] (O) -- (Pproj);
  \draw[dashed, gray!60] (Pproj) -- (P);
  \draw[gray!80, thick] (O) -- (P) node[midway, left=2pt, font=\scriptsize, fill=white, inner sep=1pt] {$\rho$};

  % Angle phi and theta arcs
  \draw[gray!80] (0,0,1.0) arc (90:50:1.0) node[midway, right=1pt, font=\tiny] {$\phi$};
  \draw[gray!80] (0.7,0,0) arc (0:35:0.7);
  \node[font=\tiny] at (0.9,0.3,0) {$\theta$};

  % Parallel latitude circle segment (showing radius r = rho*sin phi)
  \draw[dashed, quotientorange!70] (0,0,1.98) -- (P)
    node[pos=0.4, above=2pt, font=\tiny] {$r=\rho\sin\phi$};

  % Volume brick faces (shaded)
  \fill[areafill, opacity=0.4] (P) -- (Pr) -- (Prphi) -- (Pphi) -- cycle;
  \fill[areafill, opacity=0.55] (P) -- (Pr) -- (Prtheta) -- (Ptheta) -- cycle;
  \fill[areafill, opacity=0.3] (Pr) -- (Prphi) -- (Pcorner) -- (Prtheta) -- cycle;
  \fill[areafill, opacity=0.6] (Pphi) -- (Prphi) -- (Pcorner) -- (Pphitheta) -- cycle;

  % Brick outline edges
  \draw[thick, gray!50] (P) -- (Pphi);
  \draw[thick, gray!50] (P) -- (Ptheta);
  \draw[thick, gray!50] (Pphi) -- (Pphitheta);
  \draw[thick, gray!50] (Ptheta) -- (Pphitheta);
  \draw[thick, gray!50] (Pr) -- (Prtheta);
  \draw[thick, gray!50] (Prphi) -- (Pcorner);
  \draw[thick, gray!50] (Pphitheta) -- (Pcorner);
  \draw[thick, gray!50] (Prtheta) -- (Pcorner);

  % Highlighted three perpendicular differential edges
  \draw[tangentred, very thick] (P) -- (Pr)
    node[midway, left=2pt, font=\scriptsize] {$\dd\rho$};
  \draw[curveblue, very thick] (P) -- (Pphi)
    node[pos=0.8, left=2pt, font=\scriptsize] {$\rho\,\dd\phi$};
  \draw[quotientorange, very thick] (P) -- (Ptheta)
    node[pos=0.7, below=3pt, font=\scriptsize] {$\left(\rho\sin\phi\right)\dd\theta$};

  \node[font=\footnotesize, align=center] at (1.2,1.8,-1.5)
    {$\dd V = (\dd\rho)\,(\rho\,\dd\phi)\,(\rho\sin\phi\,\dd\theta) = \rho^2\sin\phi\,\dd\rho\,\dd\phi\,\dd\theta$};
\end{tikzpicture}
\caption{Geometry of the spherical volume element $\dd V$.
A tiny spherical cell is bounded by surfaces of constant $\rho$, $\phi$, and $\theta$.
Because these coordinate curves intersect at right angles, the cell forms an infinitesimal
box with radial thickness $\dd\rho$, meridional height $\rho\,\dd\phi$, and circumferential width
$(\rho\sin\phi)\,\dd\theta$ (along a circle of latitude with radius $r=\rho\sin\phi$).
Multiplying the three perpendicular dimensions yields $\dd V = \rho^2\sin\phi\,\dd\rho\,\dd\phi\,\dd\theta$.}
\label{ch34:fig-spherical}
\end{figure}

\begin{happeningbox}
Do not memorize $\rho^2\sin\phi$ -- \emph{rebuild} it. The cell is nearly a
little box with three perpendicular edges of lengths $\dd\rho$, $\rho\,\dd\phi$,
and $\rho\sin\phi\,\dd\theta$. Multiplying: $\rho\cdot\rho\sin\phi=\rho^2\sin\phi$
out front, then $\dd\rho\,\dd\phi\,\dd\theta$. Each factor is just an arc length,
``radius times angle,'' for the circle that that coordinate sweeps along. The
$\sin\phi$ appears because circles of constant latitude shrink toward the poles:
near the north pole ($\phi\approx 0$) a full lap in $\theta$ covers almost no
distance, and $\sin\phi\to 0$ records it.
\end{happeningbox}

\begin{example}[Volume of a ball, finally]
\label{ch34:ex-ball}
The ball of radius $a$ is $0\le\rho\le a$, $0\le\phi\le\pi$, $0\le\theta\le 2\pi$
-- a box in spherical coordinates. With $f=1$:
\[
  \text{Vol}=\int_0^{2\pi}\!\int_0^{\pi}\!\int_0^a
     \rho^2\sin\phi\,\dd\rho\,\dd\phi\,\dd\theta .
\]
The three integrals separate because each factor depends on a single variable:
\[
  \int_0^a\rho^2\,\dd\rho=\frac{a^3}{3},\qquad
  \int_0^\pi\sin\phi\,\dd\phi=\eval{-\cos\phi}{0}^{\pi}=2,\qquad
  \int_0^{2\pi}\dd\theta=2\pi .
\]
Multiply: $\text{Vol}=\dfrac{a^3}{3}\cdot 2\cdot 2\pi=\dfrac{4}{3}\pi a^3$. The
famous formula, derived rather than recited.
\end{example}

\section{The Jacobian: change of variables as area scaling}
\label{ch34:jacobian}

We have derived coordinate scale factors geometrically: $r$ for polar coordinates
and $\rho^2\sin\phi$ for spherical coordinates. In each case the factor arose
from the geometry of an infinitesimal coordinate cell. We now establish the
general principle governing coordinate transformations in multivariable integrals,
connecting it directly to a foundational tool from linear algebra:
the \textbf{determinant}.

\subsection{The one-variable rehearsal}

The one-dimensional analogue of this principle is integration by substitution ($u$-substitution).
When you substitute $x=g(u)$ in $\int f(x)\,\dd x$, you replace
$\dd x$ by $g'(u)\,\dd u$:
\[
  \int_a^b f(x)\,\dd x=\int_{u_a}^{u_b} f(g(u))\,g'(u)\,\dd u .
\]
The factor $g'(u)$ is a \emph{stretch factor}: near $u$, the map $g$ stretches a
little interval $\dd u$ into an interval of length $|g'(u)|\,\dd u$. (We carry an
absolute value when we ignore orientation and only care about length.) The
substitution rule \emph{is} the statement ``account for how much the
transformation stretches length.''

\begin{whybox}
The Jacobian is this same idea promoted to two and three dimensions. In one
dimension a smooth map stretches lengths by $|g'(u)|$. In two dimensions a smooth
map $(u,v)\mapsto(x,y)$ stretches \emph{areas}, and the local stretch factor is
\emph{not} a single derivative but a determinant of partial derivatives, because
near a point a smooth map looks like a \emph{linear} map, and the factor by which
a linear map scales area is exactly the absolute value of its determinant -- a
fact you proved in linear algebra. Change of variables in a multiple integral is
$u$-substitution with the word ``derivative'' replaced by the word
``determinant,'' for the simple reason that area-scaling in many dimensions is
measured by a determinant.
\end{whybox}

\subsection{Why a determinant: linearization}

Let a transformation $T$ carry the $uv$-plane to the $xy$-plane:
\[
  x=x(u,v),\qquad y=y(u,v).
\]
Take a tiny rectangle in the $uv$-plane with corner at $(u,v)$ and sides
$\dd u$ (along the $u$-axis) and $\dd v$ (along the $v$-axis). Where does $T$ send
it? Each corner moves, and because the rectangle is tiny, $T$ acts on it almost
\emph{linearly}: the image is (very nearly) a parallelogram. The side that was
$\dd u\,\vec{e}_u$ becomes the vector
\[
  \vec{a}=\left(\pdv{x}{u},\,\pdv{y}{u}\right)\dd u,
\]
because moving by $\dd u$ in the $u$-direction changes $x$ by $\pdv{x}{u}\dd u$
and $y$ by $\pdv{y}{u}\dd u$. Likewise the $\dd v$ side becomes
\[
  \vec{b}=\left(\pdv{x}{v},\,\pdv{y}{v}\right)\dd v .
\]
The image cell is the parallelogram spanned by $\vec a$ and $\vec b$, and the area
of a parallelogram spanned by two vectors is the absolute value of the
$2\times 2$ determinant of their components:
\[
  \dd A_{xy}
  =\left|\det\!\begin{pmatrix}\pdv{x}{u} & \pdv{x}{v}\\[4pt]
                              \pdv{y}{u} & \pdv{y}{v}\end{pmatrix}\right|
  \,\dd u\,\dd v .
\]

\begin{intuitionbox}
\textbf{Why the determinant measures area.}
If you have two vectors $\vec a = (a_1, a_2)$ and $\vec b = (b_1, b_2)$ in the plane, why is the area
of the parallelogram they span given by the determinant $|a_1 b_2 - a_2 b_1|$?
Enclose the parallelogram inside a rectangular bounding box of width $a_1+b_1$ and height $a_2+b_2$.
The total area of this bounding box is $(a_1+b_1)(a_2+b_2) = a_1 a_2 + a_1 b_2 + a_2 b_1 + b_1 b_2$.
To find the area of the interior parallelogram, subtract the unneeded surrounding pieces:
two corner rectangles of combined area $2(a_2 b_1)$, two right triangles along the horizontal sides of combined area $2 \cdot \tfrac12 a_1 a_2 = a_1 a_2$,
and two right triangles along the vertical sides of combined area $2 \cdot \tfrac12 b_1 b_2 = b_1 b_2$.
The subtraction leaves
\[
  (a_1+b_1)(a_2+b_2) - 2a_2 b_1 - a_1 a_2 - b_1 b_2 = a_1 b_2 - a_2 b_1 .
\]
The determinant is the exact geometric area of the parallelogram, and its absolute value accounts for orientation.
\end{intuitionbox}

\begin{figure}[H]
\centering
\begin{tikzpicture}[scale=1.0]
  \begin{scope}
    \draw[->] (-0.2,0)--(2.4,0) node[right]{\scriptsize $u$};
    \draw[->] (0,-0.2)--(0,2.4) node[above]{\scriptsize $v$};
    \fill[areafill,opacity=0.5] (0.7,0.7) rectangle (1.5,1.5);
    \draw (0.7,0.7) rectangle (1.5,1.5);
    \node at (1.1,1.1) {\scriptsize $\dd u\,\dd v$};
  \end{scope}
  \draw[->, very thick, structurepurple] (2.7,1.1) -- (3.9,1.1)
     node[midway,above]{\scriptsize $T$};
  \begin{scope}[xshift=4.4cm]
    \draw[->] (-0.2,0)--(2.6,0) node[right]{\scriptsize $x$};
    \draw[->] (0,-0.2)--(0,2.4) node[above]{\scriptsize $y$};
    \fill[areafill,opacity=0.5] (0.6,0.6) -- (1.7,0.9) -- (2.0,1.8) -- (0.9,1.5) -- cycle;
    \draw (0.6,0.6) -- (1.7,0.9) -- (2.0,1.8) -- (0.9,1.5) -- cycle;
    \draw[->,red,thick] (0.6,0.6)--(1.7,0.9) node[midway,below]{\scriptsize $\vec a$};
    \draw[->,blue,thick] (0.6,0.6)--(0.9,1.5) node[midway,left]{\scriptsize $\vec b$};
  \end{scope}
\end{tikzpicture}
\caption{Under a smooth map $T$, a tiny rectangle becomes a tiny parallelogram.
Its area is $|\det|\;\dd u\,\dd v$, where the determinant is built from the
partial derivatives. That determinant is the Jacobian.}
\label{ch34:fig-jacobian}
\end{figure}

That determinant has a name.

\begin{definition}[Jacobian]
\label{ch34:def-jacobian}
For a transformation $x=x(u,v)$, $y=y(u,v)$, the \emph{Jacobian determinant} is
\[
  \frac{\partial(x,y)}{\partial(u,v)}
  =\det\!\begin{pmatrix}\pdv{x}{u} & \pdv{x}{v}\\[4pt]
                        \pdv{y}{u} & \pdv{y}{v}\end{pmatrix}
  =\pdv{x}{u}\pdv{y}{v}-\pdv{x}{v}\pdv{y}{u}.
\]
\end{definition}

\begin{theorem}[Change of variables]
\label{ch34:thm-cov}
If $T:(u,v)\mapsto(x,y)$ is a smooth, invertible map carrying a region $S$ in the
$uv$-plane onto a region $R$ in the $xy$-plane, then
\[
  \iint_R f(x,y)\,\dd x\,\dd y
  =\iint_S f\big(x(u,v),y(u,v)\big)\,
     \left|\frac{\partial(x,y)}{\partial(u,v)}\right|\dd u\,\dd v .
\]
\end{theorem}

To change variables in a double integral: substitute the expressions for $x(u,v)$
and $y(u,v)$ into $f$, replace $\dd x\,\dd y$ by the absolute Jacobian determinant
times $\dd u\,\dd v$, and re-express the limits of integration in the new variables $(u,v)$.
This generalizes single-variable $u$-substitution to multiple dimensions.

\begin{connectionbox}
In linear algebra, composing two linear maps multiplies their matrices, and the determinant of a
product is the product of the determinants: $\det(AB) = (\det A)(\det B)$. Because smooth maps
behave locally like linear maps, their area-scaling factors multiply in the exact same way when
you chain two coordinate changes together -- the multivariable echo of $(f\circ g)' = f'\cdot g'$
in single-variable substitution.
\end{connectionbox}

\subsection{Polar coordinates, re-derived without pictures}

We can verify this formula on polar coordinates. For polar
coordinates the ``new variables'' are $(r,\theta)$ and
$x=r\cos\theta$, $y=r\sin\theta$. The partials:
\[
  \pdv{x}{r}=\cos\theta,\quad \pdv{x}{\theta}=-r\sin\theta,\quad
  \pdv{y}{r}=\sin\theta,\quad \pdv{y}{\theta}=r\cos\theta .
\]
The Jacobian:
\[
  \frac{\partial(x,y)}{\partial(r,\theta)}
  =\det\!\begin{pmatrix}\cos\theta & -r\sin\theta\\ \sin\theta & r\cos\theta\end{pmatrix}
  =r\cos^2\theta+r\sin^2\theta=r .
\]
So $\dd x\,\dd y=|r|\,\dd r\,\dd\theta=r\,\dd r\,\dd\theta$ (since $r\ge0$),
exactly the factor we found by staring at one fanned-out cell. The formula and the
picture agree, as they must. The spherical factor $\rho^2\sin\phi$ comes out of
the analogous $3\times 3$ Jacobian determinant in the same way; computing it is a
useful exercise, and the Jacobian formula applies directly even when geometric
visualization is difficult.

\begin{example}[A change of variables that is not polar]
\label{ch34:ex-cov}
Evaluate $\iint_R (x+y)\,\dd A$ where $R$ is the parallelogram with vertices
$(0,0)$, $(1,1)$, $(2,1)$, $(1,0)$ -- ugly bounds in $x,y$. Its edges follow the
lines $x-y=\text{const}$ and $y=\text{const}$, so set
\[
  u=x-y,\qquad v=y .
\]
Then $x=u+v$, $y=v$. In the new variables $R$ becomes the clean rectangle
$0\le u\le 1$, $0\le v\le 1$. The Jacobian:
\[
  \frac{\partial(x,y)}{\partial(u,v)}
  =\det\!\begin{pmatrix}\pdv{x}{u}&\pdv{x}{v}\\[3pt]\pdv{y}{u}&\pdv{y}{v}\end{pmatrix}
  =\det\!\begin{pmatrix}1&1\\0&1\end{pmatrix}=1 .
\]
The map is a shear, and shears preserve area -- determinant $1$ -- so $\dd A$ is
unchanged. With $x+y=(u+v)+v=u+2v$:
\[
  \iint_R(x+y)\,\dd A=\int_0^1\!\int_0^1 (u+2v)\cdot 1\,\dd u\,\dd v .
\]
Inner: $\int_0^1(u+2v)\dd u=\tfrac12+2v$. Outer:
$\int_0^1(\tfrac12+2v)\dd v=\tfrac12+1=\tfrac32$. The shear turned a
parallelogram into a unit square at no cost, because its Jacobian was $1$.
\end{example}

\begin{recoverybox}
\textbf{Pause and think.}
What is the Jacobian determinant of a pure diagonal scaling transformation
$x = 3u, y = 5v$?

\emph{Answer:} The matrix of partial derivatives is $\begin{pmatrix}3 & 0 \\ 0 & 5\end{pmatrix}$,
whose determinant is $3 \cdot 5 = 15$. Geometrically, the transformation stretches every
horizontal increment $\dd u$ by a factor of $3$ and every vertical increment $\dd v$ by a
factor of $5$. The area of every tiny tile scales by $3 \times 5 = 15$, so
$\dd x\,\dd y = 15\,\dd u\,\dd v$. The Jacobian measures this local area-stretching factor everywhere.
\end{recoverybox}

\begin{trapbox}
Common pitfalls when applying a change of variables:
\begin{itemize}
\item Omitting the Jacobian determinant factor entirely.
\item Forgetting the absolute value. Area is non-negative; a negative determinant
indicates that the mapping reverses orientation, so the local scaling factor is $|\det|$.
\item Forgetting to transform the region of integration: substituting expressions into $f$
and inserting the Jacobian while retaining the old Cartesian limits. Changing variables
requires re-describing the region in terms of the new coordinates $(u,v)$.
\end{itemize}
\end{trapbox}

\begin{historybox}
The idea of calculating volumes by slicing began in antiquity, but its modern
form began with Bonaventura Cavalieri in 1635. Cavalieri proposed
that two solids have identical volumes if their cross-sectional slices by parallel
planes have equal areas at every height -- treating a solid as a stack of flat
layers. Cavalieri's ``method of indivisibles'' drew criticism: critics asked how
adding infinitely many slices of zero thickness could produce a solid volume,
and the technique worked only when slicing along parallel planes.

In the eighteenth century, Leonhard Euler and Joseph-Louis Lagrange extended
slicing to iterated integrals, $\iint \dd x\,\dd y$. But physical problems
in gravitation and fluid flow required curved coordinates (polar, spherical,
and elliptical). For each new coordinate system, Euler and Lagrange had to re-derive
the volume factor by ad-hoc geometric sketches. Mathematicians also routinely
swapped the order of integration assuming it was always valid, until Cauchy and others
found counterexamples where switching orders gave different values.

The breakthrough unifying coordinate transformations came in 1841 from Carl Gustav Jacob
Jacobi. Jacobi showed that near any point, every smooth map $(u,v)\mapsto(x,y)$ behaves
like a linear transformation -- its derivative matrix -- and that the local area-scaling
factor is the absolute value of the determinant of that matrix. Six decades later, in
1907, Guido Fubini used Lebesgue integration to prove the definitive condition for
swapping integrals: the two iterated orders agree whenever $\iint |f|\,\dd A < \infty$.

In two and three dimensions, integration is well understood. But modern applications
in statistics, machine learning, and statistical physics require integrating over
hundreds or thousands of variables. Standard grid sums fail because dividing a
$100$-dimensional cube into just $10$ slices along each axis yields $10^{100}$ cells.
High-dimensional integration relies instead on randomized Monte Carlo methods and Markov
chain sampling.
\end{historybox}

\begin{aestheticbox}
The main concepts of this chapter can be summarized briefly: \emph{An integral
adds up a quantity over little cells; to compute it you slice (Fubini), and to
bend the coordinate grid you scale each cell by a determinant (the Jacobian).}
The single integral, the double, the triple; polar, cylindrical, spherical; the
$\sqrt\pi$ of the Gaussian -- all of it is that one sentence, applied with
different cells and different determinants. Calculus in many variables is not a
pile of new formulas. It is the same Riemann sum you have known since Chapter~20,
asked the same two questions: how shall I organize the sum, and how big is one
cell?

In the physical world, quantities often carry direction: wind has a velocity at each
point, and forces act along trajectories. How do you compute the work done by a force
along a curved path, or the net fluid pouring through a surface? In Chapter~35,
integration extends to vector fields, leading to line integrals, surface flux, and
Green's and Stokes' theorems relating regions to their boundaries.
\end{aestheticbox}

\section{Exercises}
\label{ch34:exercises}

\begin{exercise}[Mental check: zero by symmetry]
\label{ch34:exr-mental-sym}
Without evaluating any antiderivative, explain why
\[
  \iint_R (x^5 \cos y + y^3)\,\dd A = 0
\]
where $R$ is the rectangle $[-2,2]\times[-1,1]$. (Identify the symmetry of each term in its respective variable.)
\end{exercise}

\begin{exercise}[Mental check: predict the scaling]
\label{ch34:exr-mental-scaling}
A linear transformation from the $uv$-plane to the $xy$-plane is defined by
$x = 4u - v$ and $y = u + 2v$.
What is its Jacobian determinant? If a region $S$ in the $uv$-plane has area $3$,
what is the area of its image $R$ in the $xy$-plane?
\end{exercise}

\begin{exercise}[Warm-up over a rectangle]
\label{ch34:exr-rect}
Compute $\displaystyle\iint_R (4-x-y)\,\dd A$ over $R=[0,1]\times[0,2]$ by
iterating in both orders, and confirm the two answers agree.
\end{exercise}

\begin{exercise}[A region between curves]
\label{ch34:exr-curves}
Let $R$ be the region bounded by $y=x^2$ and $y=2x$. Sketch it, decide which
order is convenient, and evaluate $\displaystyle\iint_R x\,\dd A$. (First find
where the two curves meet.)
\end{exercise}

\begin{exercise}[Change the order]
\label{ch34:exr-order}
Rewrite
$\displaystyle\int_0^4\!\int_{\sqrt y}^{2} \frac{1}{x^3+1}\,\dd x\,\dd y$
with the order of integration reversed, then evaluate it. (As written the inner
integral is unpleasant; reversed it is a clean substitution.)
\end{exercise}

\begin{exercise}[Integration in polar coordinates]
\label{ch34:exr-polar}
Evaluate $\displaystyle\iint_R \frac{1}{1+x^2+y^2}\,\dd A$, where $R$ is the disk
$x^2+y^2\le 1$, by converting to polar coordinates. Where did the factor of $r$
make the radial integral elementary?
\end{exercise}

\begin{exercise}[A radial Gaussian]
\label{ch34:exr-gauss}
Using polar coordinates, find
$\displaystyle\iint_{\R^2} e^{-(x^2+y^2)}\,\dd A$. Then explain in one sentence
why this equals $\left(\int_{-\infty}^\infty e^{-x^2}\dd x\right)^2$, and use it
to recover $\int_{-\infty}^\infty e^{-x^2}\,\dd x=\sqrt\pi$.
\end{exercise}

\begin{exercise}[Triple integral over a box and a wedge]
\label{ch34:exr-triple}
(a) Compute $\displaystyle\iiint_E xyz\,\dd V$ over the box $[0,1]^3$.
(b) Compute the volume of the solid below $z=1-x$ and above the unit square
$[0,1]\times[0,1]$ in the $xy$-plane (assume $1-x\ge 0$ throughout).
\end{exercise}

\begin{exercise}[Spherical coordinates]
\label{ch34:exr-spherical}
Find the mass of a ball of radius $a$ whose density at a point is proportional to
its distance from the centre, $f=k\rho$. Use spherical coordinates and the volume
element $\rho^2\sin\phi\,\dd\rho\,\dd\phi\,\dd\theta$. (You should find the mass is
$k\pi a^4$.)
\end{exercise}

\begin{exercise}[Compute a Jacobian]
\label{ch34:exr-jacobian}
For the transformation $x=u^2-v^2$, $y=2uv$, compute the Jacobian
$\dfrac{\partial(x,y)}{\partial(u,v)}$ and simplify. (You should get
$4(u^2+v^2)$. This is the squaring map of complex numbers, $z\mapsto z^2$; the
area-scaling at a point grows with distance from the origin -- interpret why.)
\end{exercise}

\begin{exercise}[Change of variables, designed]
\label{ch34:exr-cov}
Evaluate $\displaystyle\iint_R e^{(x-y)/(x+y)}\,\dd A$ over the triangle with
vertices $(0,0)$, $(1,0)$, $(0,1)$ using the substitution $u=x-y$, $v=x+y$.
Find the Jacobian, convert the region, and reduce to a one-variable integral.
\end{exercise}

\begin{exercise}[Conceptual: why the $r$?]
\label{ch34:exr-polar-concept}
Without computing any integral, explain in your own words why the polar area
element carries a factor of $r$ but the Cartesian element $\dd x\,\dd y$ does not.
Then explain why $\iint_R 1\cdot\dd r\,\dd\theta$ (with the $r$ wrongly omitted)
does \emph{not} give the area of a disk, and what quantity it gives instead.
\end{exercise}

\begin{exercise}[Proof: linearity from the sum]
\label{ch34:exr-linearity}
Prove that $\iint_R\big(f+g\big)\,\dd A=\iint_R f\,\dd A+\iint_R g\,\dd A$
directly from the Riemann-sum definition, by writing out the sum for $f+g$ and
splitting it. Identify the single property of finite sums that makes the
argument work, and note that the limit step requires both integrals to exist.
\end{exercise}

\begin{exercise}[Proof: a determinant is an area scaler]
\label{ch34:exr-det-area}
Let $T$ be the linear map sending $(u,v)$ to $(au+bv,\;cu+dv)$. Show directly that
the unit square $[0,1]^2$ maps to a parallelogram of area $|ad-bc|$, and conclude
that $T$ scales every area by the factor $|ad-bc|$. Explain how this is the
constant-coefficient special case of the Jacobian theorem
(Theorem~\ref{ch34:thm-cov}).
\end{exercise}
